\appendix
\section{Anexos}
\label{sec:anexos}

\subsection{Anexo A — Especificación OpenAPI}

Incluya o referencie el archivo \texttt{openapi\_swagger/openapi.yaml}. En el repositorio mantenga también \texttt{openapi\_swagger/API.md}. Exporte la UI de Swagger (\texttt{/docs}) a PDF para el entregable gráfico.

\begin{lstlisting}[style=bashstyle]
openapi: 3.0.3
info:
  title: Microservicio ML no supervisado
  version: 1.0.0
paths:
  /api/v1/infer:
    post:
      summary: Ejecutar inferencia no supervisada
  /api/v1/inferences:
    get:
      summary: Consultar historial de inferencias
\end{lstlisting}

Sustituya el fragmento anterior por la especificación completa generada desde su framework (FastAPI exporta OpenAPI automáticamente).

\subsection{Anexo B — Enlaces de entrega}

\begin{tabular}{@{}lp{9cm}@{}}
  \toprule
  Recurso & Enlace / ruta \\
  \midrule
  Repositorio Git & \ReporteRepositorio \\
  Video demostrativo & \ReporteVideo \\
  Notebook & \pendiente{ruta en repo} \\
  Colección Postman/Bruno & \pendiente{ruta en repo} \\
  API.md & \pendiente{ruta en repo} \\
  \bottomrule
\end{tabular}

\subsection{Anexo C — Guion sugerido del video (5--10 min)}

\begin{enumerate}
  \item Contexto del proyecto móvil y problemática (1 min).
  \item Arquitectura del microservicio (1--2 min).
  \item Entrenamiento del modelo no supervisado en notebook (2 min).
  \item Demo en vivo: inferencia + consulta de historial (2--3 min).
  \item Swagger y colección de pruebas (1 min).
  \item Cómo se integrará con la app móvil al final del curso (1 min).
\end{enumerate}

\subsection{Anexo D — Registro de decisiones técnicas (opcional)}

\pendiente{Tabla breve: fecha, decisión, alternativa descartada, motivo.}
